\documentclass[../proyecto.tex]{subfiles}
\graphicspath{{\subfix{../images/}}}
\begin{document}

\section{Hipótesis y pregunta de investigación}

\subsection{Hipótesis}

La sintesíntesis, un estilo de vida que consiste en sintetizar sintetizadores, es una práctica de creación artística de herramientas electrónicas y computacionales, que cuando es realizada con los valores de ser de fuente abierta y reparables, genera nuevas autonomías electrónicas.

La sintesíntesis, un estilo de vida que consiste en sintetizar sintetizadores, es una insistencia en la luthería de instrumentos electrónicos y computacionales para emitir ruidos. Su realización de forma popular, entendida como pública, abierta, artesanal y en alternativa a las industrias, emergen nuevas estrategias de diseño y fabricación de dispositivos, aparatos, máquinas, infraestructuras y metáforas que permiten la popularización de la producción de ruidos y conocimientos técnicos.

\subsection{Pregunta de investigación}

¿Cómo el diseño y la construcción de sintetizadores sonoros, al realizarse con principios de fuente abierta, reparabilidad y legibilidad, impactan en la popularización, adopción y acceso al ruido y al desarrollo de formas de conocimiento técnico y autonomía electrónica y computacional latinoamericana?

% máximo media página

\end{document}